\chapter{Chapter 15: Resource Geometry}\label{chapter-15-resource-geometry}

\emph{In which we state the revised thesis plainly and ask what it means.}

\begin{center}\rule{0.5\linewidth}{0.5pt}\end{center}

\section{The Thesis, Revised}\label{the-thesis-revised}

After 53 findings and 63 experiments, the thesis is:

\textbf{Within the tested regime --- small arenas, locatable wells, non-monopolizable access, passive proximity benefit --- the spatial patterns we measured (clustering, incidental energy transfer, co-location, survival distributions) are determined by the geometry of resource attractors, not by the movement algorithms of agents.}

Six distinct controllers --- including agents that do not move at all --- produce all the spatial-statistical phenomena we originally attributed to a cooperation mechanism. Clustering, resonance events, co-location patterns, and survival distributions are geometric consequences of shared spatial attractors. We use the term ``social structure'' for readability but acknowledge it labels a spatial-statistical phenomenon, not a social one in the human sense.

This is by-product mutualism: selfish action produces incidental mutual benefit when the environment channels individuals toward the same locations.

\section{What This Explains}\label{what-this-explains}

Every finding from Chapters 5 through 12 is reinterpreted under resource geometry:

\textbf{Clustering} (F1-4) emerges because all agents seek the same two or three wells, not because the FEP gradient creates social structure. Any controller that navigates to wells produces identical clustering.

\textbf{Phase transitions} (F20) occur because maintenance pressure depletes energy faster than wells can regenerate it. The transition is between well-access regimes, not cooperation regimes. Below the critical pressure, well regen exceeds maintenance. Above it, the deficit grows faster than resonance can compensate.

\textbf{Social bonds} (F19) improved survival by 70\% because distance constraints held agents near wells, not near each other. The bond prevented drift away from resource access points.

\textbf{Adversarial dynamics} (F14-18) reflected competition for well access, not social conflict. Adversaries died faster because their inverted harmony profiles meant they gained no incidental resonance --- not because defection was punished, but because they received less incidental benefit from co-location.

\textbf{The equal opportunity paradox} (F26) was a race for well access, not a cooperation failure. Agents starting at the same point all competed for the nearest well, and small stochastic advantages compounded. This is resource competition, not social inequality.

\textbf{The Bowling Alone result} (F43) showed that voluntarily refusing to co-locate with other agents reduces incidental resonance. Under the revised thesis, ``bowling alone'' means voluntarily walking away from resource points where others are already gathering --- giving up the free benefit of by-product mutualism, not refusing to cooperate.

\section{What This Does Not Explain}\label{what-this-does-not-explain}

The engine cannot speak to environments where:

\begin{itemize}
\tightlist
\item
  Resources are non-spatial (information, social capital, reputation)
\item
  Resources require genuine coordination to extract (pack hunting, collective construction)
\item
  Agents have conflicting resource preferences (not all wanting the same wells)
\item
  The arena is large enough that random exploration fails
\item
  Agents can monopolize resources (excluding others from wells)
\end{itemize}

These are the conditions under which real cooperation --- not by-product mutualism --- becomes necessary. Our engine's design excludes all of them, which is why it found only by-product mutualism. This is not a flaw in the engine; it is a precise statement of the engine's scope.

\section{The Constructive Proof, Revised}\label{the-constructive-proof-revised}

Chapter 13 of the original manuscript argued that the work had value as a constructive proof of sufficiency: ``These conditions are sufficient to produce cooperation.''

The revised constructive proof is more interesting: ``These conditions are sufficient to produce \emph{the appearance} of cooperation without any cooperation mechanism.'' The social structure is real. The clustering is real. The survival patterns are real. The cooperation events are real. But they are all by-products of individual resource-seeking, not evidence of a cooperation mechanism.

This is a contribution to simulation science: it demonstrates how easily emergent behavior can be misattributed when the researcher tests the environment but not the controller.

\section{The Lesson for Simulation Science}\label{the-lesson-for-simulation-science}

We made a specific, correctable error. We varied everything except the one thing that mattered. We changed the environment 45 times, added adversaries, volatility, topology, dimensionality, learning, memory, communication --- and concluded that our findings were robust. They were robust to everything except the one test we did not run: replacing the controller.

This is not unique to our project. Any agent-based model that uses a fixed behavioral rule and varies only environmental parameters is vulnerable to the same error. The behavior may be a property of the rule, not the environment. The only way to know is to replace the rule and see if the behavior survives.

Our by-product mutualism result survived controller replacement --- cooperation appeared under all six controllers. But its \emph{mechanism} did not survive. The mechanism we attributed to the FEP gradient was actually resource geometry. The falsification was productive: it produced a cleaner, more parsimonious, and more general explanation.

We recommend that every simulation paper include a controller replacement study as standard methodology. If your finding survives controller replacement, it is environment-driven. If it does not, it is controller-driven. Both are valid findings, but they mean different things.

\section{The Closing}\label{the-closing}

We began this book expecting to demonstrate that cooperation is thermodynamically inevitable. We demonstrated instead that social structure is geometrically inevitable --- a weaker but more honest claim, and ultimately a more interesting one.

The geometry of resources determines the geometry of societies. Agents do not cooperate because they must. They co-locate because the wells are where the energy is. Everything else --- the clustering, the apparent coordination, the social structure that looks so much like community --- is free.

Physics does not cooperate. It clusters at the wells. And in that clustering, incidentally, it produces everything that looks like a society.

\begin{center}\rule{0.5\linewidth}{0.5pt}\end{center}

\emph{Next: Chapter 16 --- The Freedom to Walk Away}
